\documentclass[11pt,letterpaper]{article}
\usepackage[T1]{fontenc}
\usepackage[utf8]{inputenc}
\usepackage{lmodern}
\usepackage[margin=0.85in,headheight=15pt]{geometry}
\usepackage{amsmath,amssymb,amsthm,mathtools}
\usepackage{microtype,xcolor,booktabs,longtable,array}
\usepackage{listings,fancyhdr,enumitem,xurl}
\usepackage[colorlinks=true,linkcolor=navy,citecolor=navy,urlcolor=navy]{hyperref}
\definecolor{navy}{HTML}{17324D}
\definecolor{teal}{HTML}{176B70}
\definecolor{codebg}{HTML}{F4F6F8}
\definecolor{commentgray}{HTML}{52606D}
\hypersetup{pdftitle={Decomposing Algebraic Numbers into Low-Degree Sums and Products},
 pdfauthor={OpenAI},pdfsubject={Exact algorithms, degree optimality, and Wolfram Language implementation}}
\pagestyle{fancy}\fancyhf{}
\fancyhead[L]{\small\textsc{Low-degree algebraic decompositions}}
\fancyhead[R]{\small Exact algorithms and certificates}
\fancyfoot[C]{\thepage}
\renewcommand{\headrulewidth}{0.35pt}
\setlength{\parindent}{0pt}
\setlength{\parskip}{5.5pt plus 1pt minus 1pt}
\setcounter{tocdepth}{2}
\setlist{nosep,topsep=4pt,leftmargin=*}
\newtheorem{theorem}{Theorem}[section]
\newtheorem{proposition}[theorem]{Proposition}
\newtheorem{corollary}[theorem]{Corollary}
\newtheorem{lemma}[theorem]{Lemma}
\theoremstyle{definition}
\newtheorem{definition}[theorem]{Definition}
\newtheorem{example}[theorem]{Example}
\theoremstyle{remark}
\newtheorem{remark}[theorem]{Remark}
\newcommand{\Q}{\mathbb Q}
\newcommand{\Z}{\mathbb Z}
\newcommand{\C}{\mathbb C}
\newcommand{\Gal}{\operatorname{Gal}}
\newcommand{\Res}{\operatorname{Res}}
\newcommand{\Tr}{\operatorname{Tr}}
\newcommand{\Nm}{\operatorname{Nm}}
\newcommand{\Span}{\operatorname{span}_{\Q}}
\newcommand{\lpf}{\operatorname{lpf}}
\newcommand{\Dplus}{D_{+}}
\newcommand{\Dtimes}{D_{\times}}
\newcommand{\wl}[1]{\texttt{#1}}
\lstdefinelanguage{Wolfram}{
 morekeywords={Module,Table,Do,If,Return,Root,RootReduce,MinimalPolynomial,
 FactorList,ToNumberField,AlgebraicNumberPolynomial,PolynomialRemainder,
 Resultant,NullSpace,RowReduce,LinearSolve,MatrixRank,Function,Expand,
 Options,OptionsPattern,OptionValue,BeginPackage,Begin,End,EndPackage,
 Catch,Throw,Failure,Missing,True,False,Infinity,All,CountRoots,
 IrreduciblePolynomialQ,Get,VerificationTest,TestReport},
 sensitive=true,morecomment=[s]{(*}{*)},morestring=[b]"}
\lstset{language=Wolfram,basicstyle=\ttfamily\footnotesize,
 keywordstyle=\color{navy},commentstyle=\color{commentgray}\itshape,
 stringstyle=\color{teal},backgroundcolor=\color{codebg},
 breaklines=true,breakatwhitespace=false,columns=fullflexible,
 keepspaces=true,showstringspaces=false,upquote=true,
 frame=single,rulecolor=\color{black!12},framerule=0.35pt,
 xleftmargin=5pt,xrightmargin=5pt,framesep=5pt,
 aboveskip=8pt,belowskip=8pt,tabsize=2}
\newcommand{\scopebox}[1]{\begin{center}\fcolorbox{navy!35}{codebg}{\begin{minipage}{0.94\linewidth}#1\end{minipage}}\end{center}}
\title{\vspace{-1.2em}\color{navy}\textbf{Decomposing Algebraic Numbers\\into Low-Degree Sums and Products}\\[0.5em]
\large Exact algorithms, global degree bounds,\\and a Wolfram Language reference implementation}
\author{A solution study for Mathematica StackExchange question 105933}
\date{6 September 2026}
\begin{document}
\maketitle
\thispagestyle{empty}
\begin{abstract}
Given an exact algebraic number, we seek a flat sum or product of algebraic
numbers whose largest absolute degree is as small as possible. The two
ninth-degree examples in the question have globally optimal decompositions
into the specified pair of cubic roots. We prove their identities by explicit
polynomial congruences and prove optimality against arbitrarily many factors
or terms, including ones outside the input field. We then develop a
template-free algorithm: factor over the input number field, recover every
subfield by rational linear algebra, and solve additive membership or
multiplicative intersection problems. A normal-closure averaging argument
makes additive minimization globally complete. Products admit an exact
binary algorithm, an exact tensor test for linearly disjoint families, and a
bounded exhaustive search. We distinguish these guarantees carefully from a
general unrestricted multiplicative minimax algorithm. The accompanying
archive contains the complete Wolfram source, local regression tests, and
independently executed exact SymPy verification.
\end{abstract}
\scopebox{\textbf{What is established.}
Both examples are solved with a proved optimum of \(3\).
The additive problem has a complete finite algorithm, subject in practice to
explicit resource limits. The product implementation is complete in the
specified binary, independent-family, or bounded-search classes; it is
\emph{not} presented as a terminating solver for every unrestricted global
multiplicative minimization problem.

\textbf{Execution status.} The mathematical certificates and an independent
SymPy implementation were executed successfully. A Wolfram connector failed
with HTTP 404 before evaluation, so the supplied Wolfram package and test
suite have not been runtime-tested in Mathematica.}
\vfill
\begin{center}\small
Source files: \wl{RootDecomposition.wl}, \wl{tests.wlt}, \wl{examples.wl},\\
\wl{verify\_exact.py}, and \wl{verification\_results.json}.
\end{center}
\clearpage
{\setlength{\parskip}{0pt}\small\tableofcontents}
\clearpage

\section{The two examples: answers and an optimality certificate}
\label{sec:examples}
Write
\begin{align}
 P(x)&=x^3+x+1, & Q(x)&=x^3-x+1,\\
 f_\times(x)&=x^9+2x^7-3x^6+x^5-x^4+3x^3-x-1,\label{eq:fp}\\
 f_+(x)&=x^9+6x^6+3x^5-15x^3+24x^2-4x+8.\label{eq:fs}
\end{align}
Let \(a\) and \(b\) be the real roots of \(P\) and \(Q\), respectively.
They are the two cubic \wl{Root} objects requested in the question.
The answers are
\begin{equation}\boxed{\alpha_\times=ab,\qquad \alpha_+=a+b,\qquad
 \Dtimes(\alpha_\times)=\Dplus(\alpha_+)=3.}\label{eq:answers}\end{equation}
Here \(\alpha_\times\) and \(\alpha_+\) denote the real roots of
\eqref{eq:fp} and \eqref{eq:fs}. Their absolute degrees are both nine.
The notation \(D_+\) and \(D_\times\) is defined precisely in
Section~\ref{sec:objective}.

\subsection{Immediate Mathematica use}
The following is an exact identity check, not a numerical recognition test:
\begin{lstlisting}
a = Root[1 + # + #^3 &, 1];
b = Root[1 - # + #^3 &, 1];
ap = Root[-1 - # + 3 #^3 - #^4 + #^5 - 3 #^6 +
    2 #^7 + #^9 &, 1];
as = Root[8 - 4 # + 24 #^2 - 15 #^3 + 3 #^5 +
    6 #^6 + #^9 &, 1];

{RootReduce[ap - a b], RootReduce[as - a - b]}
(* Both entries are mathematically zero. *)
\end{lstlisting}
\wl{RootReduce} is the appropriate exact certification primitive: for
algebraic expressions built from exact roots it produces a canonical
algebraic result, with simple cases simplifying to rationals or radicals
\cite{rootreduce}. Equal numerical approximations or equal minimal
polynomials alone do not establish the desired branch identity.

To \emph{discover} the decompositions without supplying \(a\) or \(b\),
load the package accompanying this article:
\begin{lstlisting}
Get["/your/path/RootDecomposition.wl"];

rp = ProductDecomposition[ap];
rs = GlobalAdditiveDecomposition[as];

DecompositionExpression[rp]
DecompositionExpression[rs]
{rp["MaximumDegree"], rs["MaximumDegree"]}
{rp["GlobalOptimalityCertified"], rs["GlobalOptimalityCertified"]}
\end{lstlisting}
The proved mathematical targets for the last two lines are \(\{3,3\}\)
and \(\{\mathrm{True},\mathrm{True}\}\). These are expected results of
the reference implementation, not a claimed transcript from a Wolfram
kernel. Factors and terms are stored separately in the returned association;
\wl{DecompositionExpression} uses an inactive operation to preserve that
structure. The order of the two cubics is immaterial.

\subsection{Why maximum degree two is impossible}
Every algebraic number of degree at most two belongs to a quadratic extension
of \(\Q\), or to \(\Q\) itself. A field generated by any finite collection
of such numbers is multiquadratic: its degree over \(\Q\) is a power of two.
Every element of that field has degree dividing that power of two.
Consequently no number of degree nine can be a finite sum or product of
numbers of degree at most two. The displayed cubic decompositions attain
three, so three is the \emph{global} optimum. This proof places no restriction
on the number of terms or on their belonging to the input field.

Section~\ref{sec:bounds} generalizes this obstruction. For the present
examples it is already sufficient, once degree nine is certified.

\subsection{Exact recovery polynomials}
There is a particularly transparent certificate for the product. Define
\begin{align}
 A_\times(x)&=-\frac{x^7+2x^5-2x^4+1}{2},\label{eq:Ap}\\
 B_\times(x)&=-\frac{x^7+2x^5-2x^4+2x^3+2x-1}{2}.
\end{align}
Direct polynomial division over \(\Q\) gives
\begin{equation}
\begin{split}
 A_\times^3+A_\times+1&\equiv0,\\
 B_\times^3-B_\times+1&\equiv0,\\
 A_\times B_\times-x&\equiv0
\end{split}\qquad\pmod{f_\times(x)}.\label{eq:pcert}
\end{equation}
Thus the cubics are recovered from the input alone as
\(a=A_\times(\alpha_\times)\),
\(b=B_\times(\alpha_\times)\).

For the sum define
\begin{equation}
\begin{split}
 A_+(x)=-\frac1{83732}\big(&2025x^8+750x^7-4374x^6+10530x^5
 +9975x^4\\[-2pt]
 &-28868x^3-51921x^2-12496x+39992\big),
\end{split}\label{eq:As}
\end{equation}
and put \(B_+(x)=x-A_+(x)\). Then
\begin{equation}
 A_+^3+A_++1\equiv0,\qquad B_+^3-B_++1\equiv0,
 \qquad A_++B_+-x=0\pmod{f_+(x)}.\label{eq:scert}
\end{equation}
These are rational polynomial certificates, entirely independent of root
approximation. All six congruences were checked by the included Python
verification. They can also be checked locally with
\wl{PolynomialRemainder}; the complete expressions appear in
\wl{tests.wlt}.

\subsection{Degree and branch identification without numerical guessing}
The cubics are irreducible by the rational-root test. Their discriminants
are \(-31\) and \(-23\), so each has one real root and its splitting field
has Galois group \(S_3\). The two splitting fields have trivial intersection.
Indeed, a nontrivial proper common Galois quotient of two \(S_3\)-extensions
would give their unique quadratic subfields, respectively
\(\Q(\sqrt{-31})\) and \(\Q(\sqrt{-23})\), which are different; equality
of the full splitting fields would also force equality of those quadratic
subfields. Therefore the splitting fields are linearly disjoint, and the
compositum of the two chosen cubic fields has degree nine. The Galois
correspondence and compositum facts used here are standard; see
\cite[Chapters~3--4]{milne}.

The resultant identities proved in Section~\ref{sec:resultants} show that
\(ab\) and \(a+b\) satisfy \(f_\times\) and \(f_+\). Applying
\eqref{eq:pcert} or \eqref{eq:scert} at those real numbers recovers the
unique real roots of \(P\) and \(Q\). Hence each ninth-degree expression
generates the degree-nine compositum. Both displayed nonics are therefore
minimal polynomials and irreducible.

Moreover any real root \(t\) of either nonic is sent by its recovery
polynomials to real roots of \(P\) and \(Q\), which must be \(a\) and
\(b\). The last congruence then forces \(t=ab\) or \(t=a+b\), respectively.
Thus each nonic has exactly one real root. Mathematica orders the real roots
first, increasingly \cite{root}; the index \(1\) in both inputs is therefore
correct. The independent verification also checks irreducibility and exact
real-root counts directly.

\section{What is being minimized?}\label{sec:objective}
\begin{definition}
For an algebraic number \(\gamma\), its \emph{absolute degree} is
\[
 D(\gamma)=[\Q(\gamma):\Q].
\]
For an algebraic \(\alpha\), define
\[
\Dplus(\alpha)=\min_{\substack{r\geq1\\\alpha=\beta_1+\cdots+\beta_r}}
 \max_i D(\beta_i),
\qquad
\Dtimes(\alpha)=\min_{\substack{r\geq1\\\alpha=\beta_1\cdots\beta_r}}
 \max_i D(\beta_i).
\]
The terms and factors range over \(\overline{\Q}\). For nonzero \(\alpha\)
the factors are necessarily nonzero. We use \(\{0\}\) as the trivial product
representation of zero. Rationals, including zero, have degree one.
\end{definition}
The singleton representation \(\{\alpha\}\) always exists, so both minima
exist and lie between one and \(D(\alpha)\). Requiring at least two entries
would not change them: one can append zero to a sum or one to a product.
This is a minimax problem, not a minimization of expression length.

The coefficients of each defining polynomial must be rational. Otherwise
one could put \(x-\alpha\) into a degree-one \wl{Root}, making the question
vacuous. Likewise a nested radical with low \emph{relative} degrees is not
necessarily a small-absolute-degree decomposition. For example, a square
root of a quadratic irrational can have degree four over \(\Q\).

We distinguish three domains of search. An \emph{internal} decomposition
requires every term or factor to belong to \(K=\Q(\alpha)\). An
\emph{ambient-field} decomposition permits a specified finite extension
\(K\) containing \(\alpha\). A \emph{global} decomposition permits all
algebraic numbers. These are not interchangeable: Section~\ref{sec:external}
gives explicit examples with internal optimum six and global optimum four.

For an exact \wl{Root}, use \wl{MinimalPolynomial} before measuring degree.
The polynomial originally written inside a root expression need not be
irreducible. Mathematica's primitive integral minimal-polynomial convention
also supplies a consistent coefficient-height measure for the bounded
search \cite{minpoly}.

\section{Universal lower bounds}\label{sec:bounds}
\begin{theorem}[Largest-prime-factor bound]\label{thm:lpf}
Let \(n=D(\alpha)\), and put \(\lpf(1)=1\). Then
\[
 \Dplus(\alpha)\geq\lpf(n),\qquad
 \Dtimes(\alpha)\geq\lpf(n).
\]
More generally, if \(\alpha\) belongs to a field generated by finitely many
algebraic numbers of degree at most \(d\), every prime divisor of
\(D(\alpha)\) is at most \(d\).
\end{theorem}
\begin{proof}
Let the generators have degrees \(m_1,\ldots,m_r\leq d\), and let \(N_i\)
be their normal closures. The group \(\Gal(N_i/\Q)\) embeds into
\(S_{m_i}\), so its order divides \(m_i!\). The Galois group of the
compositum \(N=N_1\cdots N_r\) embeds into the product of those groups.
Therefore every prime divisor of \([N:\Q]\) is at most \(d\).
Since \(\Q(\alpha)\subseteq N\), the tower formula gives
\(D(\alpha)\mid[N:\Q]\). Any sum or product of the generators belongs
to that compositum, proving the assertions.
\end{proof}
For prime degree \(n\), neither operation can lower the maximum degree at
all. For degree nine the bound is three; for degree eight it is two. The
bound may be weaker than the true optimum, but it is inexpensive and remains
valid with arbitrarily many factors outside the input field.

A tempting stronger statement, \(D(\alpha)\mid\prod_i D(\beta_i)\), is
\emph{false} in general. Two distinct conjugate roots of an irreducible
\(S_3\)-cubic can generate its degree-six splitting field although the
product of their individual degrees is nine. The safe general statements
are the normal-closure divisibility used above and
\([\Q(\beta_1,\ldots,\beta_r):\Q]\leq\prod_iD(\beta_i)\).

\begin{proposition}[Galois-exponent obstruction]\label{prop:exponent}
Let \(G\) be the Galois group of the normal closure of \(\Q(\alpha)\).
If \(\alpha\) is a sum or product of numbers of degree at most \(d\), then
\[
 \exp(G)\mid\operatorname{lcm}(1,2,\ldots,d).
\]
\end{proposition}
\begin{proof}
The exponent of \(S_m\) is \(\operatorname{lcm}(1,\ldots,m)\), because
permutation orders are least common multiples of cycle lengths. In the
notation of the previous proof, the exponent of \(\Gal(N/\Q)\), a
subgroup of a product of such symmetric groups, divides
\(\operatorname{lcm}(1,\ldots,d)\). The normal closure of
\(\Q(\alpha)\) is a Galois subfield of \(N\), so \(G\) is a quotient
of \(\Gal(N/\Q)\). Quotients cannot increase the exponent.
\end{proof}
In particular an \(S_4\) normal closure rules out degree at most three,
because \(S_4\) contains an element of order four, whereas
\(\operatorname{lcm}(1,2,3)=6\). This sharper obstruction will certify the
sextic counterexamples later. The package computes the inexpensive
largest-prime bound automatically, not arbitrary Galois groups.

\section{Resultants: forward construction and inverse templates}
\label{sec:resultants}
\subsection{The two basic identities}
Suppose \(P,Q\in\Q[x]\) are monic, of degrees \(r,s\), with roots
\(a_i\) and \(b_j\). Then
\begin{align}
 R_+(x)&=\Res_y\bigl(P(y),Q(x-y)\bigr)
       =\prod_{i=1}^{r}\prod_{j=1}^{s}(x-a_i-b_j),\label{eq:ressum}\\
 R_\times(x)&=\Res_y\bigl(P(y),y^sQ(x/y)\bigr)
       =\prod_{i=1}^{r}\prod_{j=1}^{s}(x-a_i b_j).\label{eq:resprod}
\end{align}
To prove them, apply the defining product formula for the resultant of a
monic polynomial:
\[\Res_y(P(y),H(y))=\prod_i H(a_i).\]
In \eqref{eq:resprod}, the second argument is the polynomial
\(\sum_{j=0}^{s}q_jx^jy^{s-j}\); no division by a potentially zero root
is required. Wolfram Language supplies \wl{Resultant} directly
\cite{resultant}.

A compact implementation is
\begin{lstlisting}
RootSumPolynomial[p_, q_, x_Symbol] := Module[{y = Unique["y"]},
  Expand[Resultant[p /. x -> y, q /. x -> x - y, y]]
];
RootProductPolynomial[p_, q_, x_Symbol] := Module[
  {y = Unique["y"], s = Exponent[q, x], reversed},
  reversed = Sum[Coefficient[q, x, j] x^j y^(s-j), {j, 0, s}];
  Expand[Resultant[p /. x -> y, reversed, y]]
];
\end{lstlisting}
Substituting \(P=x^3+x+1\) and \(Q=x^3-x+1\) gives exactly
\eqref{eq:fs} and \eqref{eq:fp}. These resultant calculations were checked
independently with exact rational arithmetic.

The companion-matrix viewpoint is equivalent: the eigenvalues of
\(C_P\otimes I_s+I_r\otimes C_Q\) are all sums, and those of
\(C_P\otimes C_Q\) are all products. The original discussion contains
inverse coefficient-matching approaches using companion matrices
(yarchik) and resultants (Daniel Lichtblau) \cite{sequestion}.
Our subfield algorithms below do not require guessing those coefficient
templates.

\subsection{What coefficient matching does and does not prove}
One can choose degrees \(r,s\), introduce rational unknown coefficients in
\(P,Q\), and equate the resultant to the input's minimal polynomial. This
is often effective when the intended polynomials have very small
coefficients. There are, however, several separate requirements.

First, in general one only has
\[
 f_\alpha(x)\mid R_+(x)\quad\hbox{or}\quad f_\alpha(x)\mid R_\times(x).
\]
Equality follows when the monic polynomials have the same degree. Otherwise
there may be several irreducible factors or repeated roots. For instance,
using \(P=Q=x^2-2\) in the sum construction gives
\(x^2(x^2-8)\); using \(P=x^2-2\), \(Q=x^2-3\) in the product
construction gives \((x^2-6)^2\). The chosen sum or product has a minimal
polynomial of smaller degree.

Second, a polynomial system solved over \(\C\) can return algebraic
coefficients for \(P\) or \(Q\). Such an answer need not meet the desired
\emph{absolute} degree bounds. Rationality of the defining coefficients
must be enforced or checked. This is not fixed by merely using a
symbolic solver instead of numerical fitting.

Third, the correct root indices still have to be selected. Enumerate roots
of the recovered rational polynomials and retain combinations for which
\wl{RootReduce[target - expression]} is zero. The resultant encodes all
conjugate pairs, not the particular input branch.

Finally, normalization must preserve the rational search space. For sums,
replacing \((a,b)\) with \((a-c,b+c)\), \(c\in\Q\), allows the trace of
one chosen polynomial to be normalized to zero. For products, rescaling
\((a,b)\) to \((ca,b/c)\) is safe for rational \(c\), but forcing an
arbitrary constant coefficient to become one can require an irrational
\(r\)-th root. That restriction would exclude valid rational templates.

Thus inverse resultants are useful candidate generators and verification
tools, but solving one template is not a proof of global minimax optimality.
Factoring an irreducible input polynomial over \(\Q\) is not a substitute:
it factors a polynomial, not the selected algebraic number as a product.

\section{Recovering every subfield by exact linear algebra}
\label{sec:subfields}
The important change of viewpoint is to work in a number field rather than
to guess small polynomial coefficients. Let
\[
 K=\Q(\theta),\qquad n=[K:\Q],\qquad f(\theta)=0,
\]
where \(f\) is the monic minimal polynomial. Elements of \(K\) are represented
by rational coordinate vectors in the power basis
\(1,\theta,\ldots,\theta^{n-1}\). A subfield is represented by a matrix
whose rows are a rational basis of that subfield.

The construction uses \emph{principal subfields}. The fact that all
intermediate fields arise as intersections of principal ones, and the use
of factorization followed by linear equations, belong to the established
subfield literature. Szutkoski and van Hoeij give the formulation and faster
partition-based methods \cite{szutkoski}. We use a direct rational-matrix
version to keep the implementation self-contained; we do not claim their
optimized partition algorithm or its performance bounds for this package.

\subsection{Principal fields and completeness}
Factor \(f\) into monic irreducibles over \(K\):
\[
 f(x)=g_1(x)\cdots g_t(x),\qquad g_1(x)=x-\theta.
\]
For each factor define
\begin{equation}
 L_i=\{h(\theta):h\in\Q[x],\ \deg h<n,
             \quad h(x)-h(\theta)\equiv0\pmod{g_i(x)}\}.
 \label{eq:principal}
\end{equation}
\begin{proposition}\label{prop:principal}
Each \(L_i\) is a subfield of \(K\). Every intermediate field
\(\Q\subseteq F\subseteq K\) is an intersection of a subset of the \(L_i\).
\end{proposition}
\begin{proof}
Let \(N\) be a normal closure of \(K\). For each root \(\theta'\) of
\(g_i\), the map sending \(\theta\) to \(\theta'\) defines a
\(\Q\)-embedding \(\sigma:K\to N\). Because \(g_i\) is separable,
condition \eqref{eq:principal} says exactly that all these embeddings fix
\(h(\theta)\). The set of elements fixed by a collection of embeddings
is closed under addition, subtraction, multiplication, and inversion of
nonzero elements. It contains \(\Q\), proving the first assertion.

For an intermediate \(F\), let \(f_F\) be the minimal polynomial of
\(\theta\) over \(F\). Viewed in \(K[x]\), it is the product of a subset
of the factors \(g_i\). An element of \(F\) is fixed by every
\(F\)-embedding of \(K\), so it belongs to the intersection associated
to that subset. Conversely, an element in this intersection is fixed by all
those embeddings. Every automorphism of \(N\) fixing \(F\) restricts to
one of them; the Galois correspondence then puts the element in \(F\).
This proves equality.
\end{proof}
In particular, the field belonging to \(x-\theta\) is all of \(K\).
We need not construct \(N\) to execute this proof's algorithm: the
factorization and the linear equations take place entirely in \(K\).

\subsection{The actual rational matrix}
For one factor \(g\) of degree \(m\), compute
\[
 r_j(x)=\operatorname{rem}_{g(x)}(x^j-\theta^j),
 \qquad 0\leq j<n.
\]
Each \(r_j\) has \(m\) coefficients in \(K\), padded with zeros as
needed. Replace every coefficient by its \(n\)-entry rational coordinate
vector. Concatenation produces a vector \(c_j\in\Q^{mn}\). Then
\begin{equation}
 [\,c_0\ c_1\ \cdots\ c_{n-1}\,]v=0
 \label{eq:principalmatrix}
\end{equation}
is precisely the condition that the element with coordinates \(v\)
belong to \(L_g\). A rational nullspace therefore gives its basis.
There is no search over coefficient heights in this step.

To intersect two row spaces \(U,V\subseteq\Q^n\), use
\begin{equation}
 U\cap V=(U^\perp+V^\perp)^\perp.\label{eq:intersection}
\end{equation}
Each orthogonal complement is a rational nullspace. Store every result in
reduced row-echelon form with zero rows deleted; this is a canonical
representation, so duplicate fields can be recognized by exact equality.
Starting with \(K\), process the principal fields one at a time and retain
both every current field and its intersection with the new one. By
Proposition~\ref{prop:principal}, the final collection is the complete
subfield lattice. Explicitly forming the normal-closure Galois group is
unnecessary.

\subsection{Wolfram coordinate discipline}
Use \wl{FactorList} with the option \wl{Extension -> theta} to factor over
the ambient field. Coordinate conversion uses \wl{ToNumberField} and
\wl{AlgebraicNumberPolynomial}; reduction uses \wl{PolynomialRemainder}.
The remaining computations use exact matrix routines
\cite{factorlist,tonumberfield,anpoly,polyrem,linearsolve}.
One implementation detail is easy to get wrong. For a nonintegral primitive
element, \wl{ToNumberField} may change the generator. Coordinates relative
to that new generator cannot be used as coordinates relative to the old one.

The package avoids this by choosing an integral primitive element. If the
primitive integral minimal polynomial of the original input is
\(cx^n+a_{n-1}x^{n-1}+\cdots+a_0\), take \(\theta=c\alpha\). Substituting
\(\alpha=\theta/c\) and multiplying by \(c^{n-1}\) produces a monic
integer polynomial for \(\theta\), and \(\Q(\theta)=\Q(\alpha)\).
For integral \(\theta\), the documented representation preserves
\(\theta\) \cite{tonumberfield}. The code nevertheless checks the returned
generator before accepting coordinates. Rationals are handled separately,
including the degree-one ambient field, where one must avoid evaluating
\(0^0\) in a power basis.

For the two nonic examples, exact factorization over their own fields has
factor degrees
\[
 \{1,2,2,4\},
\]
and the recovered subfield degrees are
\begin{equation}\{1,3,3,9\}.\label{eq:examplefields}\end{equation}
The two cubic fields are \(\Q(a)\) and \(\Q(b)\). The independent
implementation reproduced these results and checked that each reported
basis contains one and is closed under multiplication of basis elements.

There is a useful structural interpretation. Over \(\Q(a)\), the three
conjugates \(a+b_j\) have a cubic whose quadratic coefficient is \(-3a\),
so its trace recovers \(a\). The three products \(ab_j\) satisfy
\(x^3-a^2x+a^3\), whose coefficients also recover \(a\), since
\(a=-a^3-1\). The full matrix algorithm generalizes this observation
without requiring us to recognize the blocks or the desired cubics first.

\section{A complete additive algorithm}
\label{sec:additive}
\subsection{All numbers in one finite ambient field}
Fix a finite field \(K\) containing \(\alpha\), and let
\(\mathcal F_d(K)\) be the set of its subfields of degree at most \(d\)
over \(\Q\). Define the rational vector space
\[
 V_d(K)=\sum_{E\in\mathcal F_d(K)} E\subseteq K.
\]
The plus sign here is a vector-space sum; it is not a field compositum.

\begin{theorem}[Additive membership criterion]\label{thm:additive}
There is a decomposition \(\alpha=\sum_i\beta_i\), with
\(\beta_i\in K\) and \(D(\beta_i)\leq d\), if and only if
\(\alpha\in V_d(K)\).
\end{theorem}
\begin{proof}
For a proposed decomposition, each \(\Q(\beta_i)\) is one of the fields
in \(\mathcal F_d(K)\), proving necessity. Conversely, expressing
\(\alpha\) in the vector-space sum gives one element of each of finitely
many subfields. Each element's degree is at most its field's degree.
\end{proof}
Concatenate row bases for all fields of degree at most \(d\), transpose,
and solve the rational linear system for the coordinate vector of
\(\alpha\). If it is inconsistent, there is no such sum in \(K\),
regardless of how many terms are allowed. If it is consistent, group the
solution coefficients according to their subfield to obtain the terms.
Under-determined systems are harmless: any exact particular solution is
sufficient \cite{linearsolve}.

Now increase \(d\) through the distinct subfield degrees. Membership changes
only at those thresholds. The first successful threshold is the true
minimum for arbitrary finite sums in \(K\). This is not a greedy recursive
splitting algorithm, and it does not assume that an optimal sum has two terms.

\subsection{Presentation and number of terms}
The linear system may return shifted or scaled generators rather than the
prettiest roots. A safe additive cleanup is trace centering. If \(\beta\)
has primitive minimal polynomial
\(c_dx^d+c_{d-1}x^{d-1}+\cdots\), then its mean conjugate is the rational
number \(-c_{d-1}/(dc_d)\). Subtract that rational from the term, and
collect all subtracted means into one rational offset. This preserves the
sum and cannot increase a term's degree. In the supplied example it removes
the rational shifts and yields the two trace-zero cubic roots.

Minimizing the number of nonzero terms is a different optimization. After
fixing the best degree \(d\), one could enumerate subsets of
\(\mathcal F_d(K)\) to find the smallest collection whose vector-space
sum contains \(\alpha\). That finite combinatorial search is not performed
by the package.

For example,
\begin{equation}
 \eta=\sqrt2+\sqrt3+\sqrt6,\qquad
 f_\eta(x)=x^4-22x^2-48x-23,\label{eq:threeadd}
\end{equation}
has optimal degree two, with three natural quadratic terms. Inside the
biquadratic field \(\Q(\sqrt2,\sqrt3)\), its three nonrational character
components cannot be supplied by only two quadratic subfields. Nor can two
external quadratic fields help: if two quadratic terms produced the
primitive degree-four \(\eta\), their compositum would have degree four
and equal \(\Q(\eta)\), putting them back in that same field. This is
an explicit reason not to substitute a binary additive search for
Theorem~\ref{thm:additive}.

\subsection{Descent to a normal closure}
The internal field \(\Q(\alpha)\) is not always enough. For sums there
is nevertheless a finite ambient field that is always enough.

\begin{theorem}[Normal-closure descent for sums]\label{thm:descent}
Let \(L/\Q\) be finite Galois and contain \(\alpha\). If
\(\alpha=\sum_{i=1}^r\beta_i\) with algebraic \(\beta_i\), then there
are \(\gamma_i\in L\) satisfying
\[
 \alpha=\sum_{i=1}^r\gamma_i,\qquad
 \gamma_i\in L\cap\Q(\beta_i),\qquad
 D(\gamma_i)\leq D(\beta_i).
\]
In particular, global additive minimization can be performed in the normal
closure of \(\Q(\alpha)\).
\end{theorem}
\begin{proof}
Choose a finite Galois extension \(M/\Q\) containing \(L\) and all
\(\beta_i\), and write
\(G=\Gal(M/\Q)\), \(H=\Gal(M/L)\). Since \(L/\Q\) is Galois,
\(H\) is normal in \(G\). Define the averaging map
\[
 \mathcal P(v)=\frac1{|H|}\sum_{h\in H}h(v)
             =\frac{\Tr_{M/L}(v)}{[M:L]}.
\]
Its values lie in \(L\), it fixes \(L\) pointwise, and it is
\(\Q\)-linear. Normality of \(H\) gives, for \(g\in G\),
\[
 g\mathcal P(v)=\frac1{|H|}\sum_{h\in H}(ghg^{-1})g(v)
              =\mathcal P(gv).
\]
Thus \(\mathcal P\) commutes with every \(g\in G\).

Set \(J_i=\Gal(M/\Q(\beta_i))\). Every member of \(J_i\) fixes
\(\beta_i\), so equivariance implies that it fixes
\(\gamma_i=\mathcal P(\beta_i)\). The Galois correspondence yields
\(\gamma_i\in\Q(\beta_i)\), while the construction already gives
\(\gamma_i\in L\). Hence its degree is no larger. Finally,
\(\sum_i\gamma_i=\mathcal P(\alpha)=\alpha\). Zero terms may be omitted.
\end{proof}
The normality assumption is essential in the equivariance step. Averaging
over \(\Gal(M/\Q(\alpha))\) without normality does not give the same
degree-preservation conclusion.

\subsection{Constructing and optimizing in that field}
An exact global procedure is now finite: construct the normal closure,
compute its finitely many subfields, and apply the linear membership test.
In Wolfram Language, a common primitive element for all roots of the
minimal polynomial can be obtained using
\begin{lstlisting}
roots = Table[Root[Function[Evaluate[f /. x -> Slot[1]]], j],
  {j, Exponent[f, x]}];
common = ToNumberField[roots, All];
theta = common[[1, 1]];
\end{lstlisting}
The \wl{All} argument matters: the documentation guarantees the smallest
common extension with \wl{All}, while the default can use a larger one
\cite{tonumberfield}. The number field generated by all these roots is
the required splitting field.

The package first tries the smaller input field. It avoids the normal
closure when the candidate meets the universal lower bound, or when the
input field is already Galois.

Otherwise the global routine constructs the closure and runs the complete
additive test there. Subject to unlimited computational resources, this is
an algorithm for the fully unrestricted additive question, not merely a
bounded coefficient search. The normal closure can have degree as large as
\(n!\), because its Galois group embeds into \(S_n\). Completeness should
therefore not be confused with uniform practical efficiency.

\section{Products of two factors: a linear, not nonlinear, test}
\label{sec:binary}
Fix nonzero \(\alpha\in K\) and two subfields \(E,F\subseteq K\).
The equation \(\alpha=uv\) initially appears bilinear in the unknown
coordinates of \(u\) and \(v\). Inversion in a field makes it linear.

\begin{theorem}[Binary multiplicative intersection]\label{thm:binary}
For \(\alpha\ne0\),
\[
 \alpha\in E^\times F^\times
 \quad\Longleftrightarrow\quad
 (\alpha E)\cap F\ne\{0\}.
\]
If \(w\in E\setminus\{0\}\) and \(\alpha w\in F\), then
\(\alpha=w^{-1}(\alpha w)\) is the required factorization.
\end{theorem}
\begin{proof}
If \(\alpha=uv\), choose \(w=u^{-1}\in E\); then \(\alpha w=v\in F\)
and is nonzero. The converse is the displayed factorization, using closure
of \(E\) under inversion.
\end{proof}
Let \(e_1,\ldots,e_r\), \(f_1,\ldots,f_s\) be rational bases. Solve
\begin{equation}
 \sum_{i=1}^r c_i\alpha e_i-\sum_{j=1}^s d_jf_j=0,
 \qquad c_i,d_j\in\Q.\label{eq:binarymatrix}
\end{equation}
Any nonzero nullspace vector gives a nonzero
\(w=\sum_i c_ie_i\). Indeed, if \(w=0\), independence of the \(e_i\)
forces all \(c_i=0\), and independence of the \(f_j\) then forces all
\(d_j=0\), a contradiction. Thus no nonlinear rational-point solver or
numerical optimization is needed for this binary problem.

Using a complete subfield lattice, enumerate all unordered pairs \(E,F\)
in increasing order of \(\max([E:\Q],[F:\Q])\). This finds the exact
minimum for \emph{two factors in \(K\)}. Every candidate is finally
converted with \wl{RootReduce}, and its actual degree and exact product are
checked. A rational sign adjustment makes odd-degree factors have positive constant
coefficient. The implementation also normalizes a monic odd-degree constant
coefficient to one when, and only when, its required root is rational.
Compensating in the last factor preserves the product. This rationality
check is essential: an arbitrary algebraic rescaling is not permitted.

For the degree-nine product, the two cubic subfields make
\eqref{eq:binarymatrix} singular, yielding the recovery factors of
Section~\ref{sec:examples}. The degree-three answer meets
Theorem~\ref{thm:lpf}, so this restricted search happens to certify a
fully global optimum.

A binary optimum alone is not sufficient for arbitrary flat products.
For example, every product of two quadratic numbers has degree at most
four, while Section~\ref{sec:tensor} exhibits a degree-eight number that is
a product of three quadratic numbers. Also, a binary optimum in
\(\Q(\alpha)\) can exceed the global binary optimum; see
Section~\ref{sec:external}.

\section{Several factors and the exact tensor test}\label{sec:tensor}
\subsection{Independent field families}
Suppose \(E_1,\ldots,E_r\subseteq K\) have degrees
\(d_1,\ldots,d_r\), and the multiplication map
\begin{equation}
 E_1\otimes_\Q\cdots\otimes_\Q E_r\longrightarrow K,
 \quad u_1\otimes\cdots\otimes u_r\longmapsto u_1\cdots u_r
 \label{eq:multtensor}
\end{equation}
is an isomorphism. Equivalently, their product basis spans \(K\) and
\(\prod_i d_i=[K:\Q]\). This is the precise independence condition
needed here; pairwise trivial intersections alone are not used as a
substitute.

Choose bases \(e_{i,1},\ldots,e_{i,d_i}\). Express \(\alpha\) uniquely as
\[
 \alpha=\sum_{j_1,\ldots,j_r}c_{j_1\cdots j_r}
       e_{1,j_1}\cdots e_{r,j_r}.
\]
Then \(\alpha\) is a product \(u_1\cdots u_r\), \(u_i\in E_i\), if
and only if the rational coefficient tensor \(c\) is a pure tensor, also
called a rank-one tensor. This is an equivalence because
\eqref{eq:multtensor} is injective as well as surjective.

\begin{proposition}[Pivot certificate for a pure tensor]\label{prop:tensor}
Let \(c\ne0\), and choose a multi-index \(p=(p_1,\ldots,p_r)\) with
\(c_p\ne0\). Put
\[
 v_{i,j}=\frac{c_{p_1,\ldots,p_{i-1},j,p_{i+1},\ldots,p_r}}{c_p}.
\]
The tensor is rank one if and only if every entry satisfies
\begin{equation}
 c_{j_1\cdots j_r}=c_p\prod_{i=1}^r v_{i,j_i}.\label{eq:pivottest}
\end{equation}
When it does, take \(u_i=\sum_jv_{i,j}e_{i,j}\), absorbing \(c_p\)
into the first factor.
\end{proposition}
\begin{proof}
If \(c\) is a product of vectors, a nonzero pivot means each selected
vector coordinate is nonzero. Dividing the appropriate slice by the pivot
normalizes that vector's selected coordinate to one, which gives the formula.
Conversely, \eqref{eq:pivottest} explicitly writes \(c\) as a product of
vectors and hence supplies the factors.
\end{proof}
The test is exact and inexpensive after the basis change. One should not
replace it by a floating-point tensor-rank estimate. For more than two
factors, general tensor rank is a different problem; here only rank one is
being tested, with the complete entrywise certificate above.

\subsection{A three-factor example that defeats binary minimization}
Consider
\begin{equation}
 \xi=(1+\sqrt2)(1+\sqrt3)(1+\sqrt5).\label{eq:threeprod}
\end{equation}
Its minimal polynomial, verified exactly, is
\begin{equation}
\begin{split}
 f_\xi(x)={}&x^8-8x^7-256x^6-768x^5+3008x^4\\
            &+6144x^3-16384x^2+4096x+4096.
\end{split}\label{eq:f8}
\end{equation}
The multiquadratic field \(\Q(\sqrt2,\sqrt3,\sqrt5)\) has degree eight.
Its standard eight radical monomials are linearly independent, and
\eqref{eq:threeprod} has a nonzero coefficient on every one. No nonidentity
sign-change automorphism fixes it, so it generates the full field. This also
proves degree eight independently of the polynomial computation.

The three natural quadratic fields satisfy \eqref{eq:multtensor}; the
coefficient tensor is the product of the three vectors \((1,1)\).
Consequently \(\Dtimes(\xi)=2\). A product of two quadratic factors is
impossible, since their compositum has degree at most four. The full field
has sixteen subfields: seven quadratic, seven quartic, \(\Q\), and itself,
as follows from the subspaces of its elementary abelian Galois group.
The independent implementation recovered all sixteen.

\subsection{Exactly what the package searches}
At each degree threshold, \wl{ProductDecomposition} first runs the binary
intersection tests. It then considers families of at least three distinct
proper nonrational subfields, tests whether their degrees multiply to
\([K:\Q]\), checks that their product basis has full rank, and applies
\eqref{eq:pivottest}. A family needs at most
\(\lfloor\log_2[K:\Q]\rfloor\) members. Repeating the same subfield is
unnecessary because factors in one field can be multiplied together.

With complete family enumeration, the result is a true minimum over the
union of the binary and independent-family classes in \(K=\Q(\alpha)\).
It does not cover arbitrary overlapping families of three or more fields,
nor factors outside \(K\). The option \wl{"MaxTensorFamilies"} limits
candidate families; if it truncates a relevant search, the result explicitly
loses the corresponding restricted-optimality claim. Meeting the universal
lower bound still certifies a global optimum independently of that
truncation.

\section{Why restricting to the input field can miss the optimum}
\label{sec:external}
Consider the quartic
\[
 h(t)=t^4-t-1
\]
and let \(r_1,r_2\) be its two real roots. Define
\(s=r_1+r_2\) and \(u=r_1r_2\). Then
\begin{equation}
 f_s(x)=x^6+4x^2-1,
 \qquad
 f_u(x)=x^6+x^4-x^3-x^2-1.\label{eq:sextics}
\end{equation}
These sextics provide counterexamples for \emph{both} operations:
\begin{equation}
 D_+^{\Q(s)}(s)=6>4=\Dplus(s),\qquad
 D_\times^{\Q(u)}(u)=6>4=\Dtimes(u).\label{eq:counteropt}
\end{equation}
The superscript denotes the restriction that every term or factor belong
to the indicated field, with arbitrarily many entries allowed.

\subsection{Deriving the sextics}
Let the remaining roots be \(r_3,r_4\), and set \(v=r_3r_4\).
Because the total root sum is zero, factorization of \(h\) gives
\[
 h(t)=(t^2-st+u)(t^2+st+v).
\]
Comparison of coefficients yields
\[
 u+v=s^2,\qquad s(u-v)=-1,\qquad uv=-1.
\]
Squaring the middle equation and using the other two gives
\(s^2(s^4+4)=1\), proving the first polynomial in \eqref{eq:sextics}.
Substitute \(v=-1/u\) and \(s^2=u-1/u\) into the squared equation to get
\((u-1/u)(u+1/u)^2=1\); multiplication by \(u^3\) gives the second.

\subsection{Galois group and field lattice}
Modulo two, \(h\) is irreducible: it has no linear factor and is not the
square of the only irreducible quadratic \(t^2+t+1\). Modulo seven, its
factorization has degrees one and three. These squarefree reductions imply
a four-cycle and a three-cycle in the Galois group, respectively, by the
standard factorization-cycle correspondence \cite[Chapter~4]{milne}.
Thus its group has order divisible by twelve. It cannot be the index-two
subgroup \(A_4\), which has no four-cycle, so the group is \(S_4\).
Its discriminant is \(-283\), consistently with these reductions.

The six unordered pairs of roots have distinct sums and distinct products.
Two overlapping pairs with the same sum or product would force equal
roots (the roots are nonzero). Equality for complementary pairs would give
\(s=0\) in the sum case or \(u=v\) in the product case; either contradicts
\(s(u-v)=-1\). Hence the stabilizer of each pair sum or pair product is
exactly
\[
 H=\langle(12),(34)\rangle\simeq S_2\times S_2,
\]
which has order four and index six. This proves that both sextics in
\eqref{eq:sextics} are minimal polynomials.

Subfields of either degree-six field correspond to overgroups of \(H\)
in \(S_4\). Apart from \(H\) and \(S_4\), the only such overgroup is
its order-eight normalizer \(D_8\). To see this, a proper intermediate
overgroup can have order eight or twelve. An order-twelve group would be
\(A_4\) and cannot contain the transpositions in \(H\). In an order-eight
overgroup, \(H\) has index two and is normal, so that overgroup is contained
in its order-eight normalizer and must equal it.

Consequently each sextic field has precisely the subfield degrees
\(1,3,6\), with a unique proper nontrivial cubic subfield. Every element
of degree below six in it belongs to that cubic. Any finite sum or product
of such elements remains in the cubic and cannot equal a primitive sextic
element. This proves the internal optima in \eqref{eq:counteropt}.

On the other hand, \(s=r_1+r_2\) and \(u=r_1r_2\) already use only
quartic numbers. The core of \(H\) in \(S_4\) is trivial, so the normal
closure of either sextic field has group \(S_4\). Proposition
\ref{prop:exponent} rules out maximum degree at most three. The global
optima are therefore exactly four.

\subsection{Practical consequences}
The subfield lattice of \(\Q(\alpha)\) is a powerful first search space,
not an automatic global completeness certificate. In the sextic sum case,
\wl{AdditiveDecomposition} correctly reports internal degree six, whereas
\wl{GlobalAdditiveDecomposition} must pass to the degree-24 normal closure
to attain the global degree four. The latter, expensive test is included
as an optional example rather than described as executed Wolfram output.

There is also no justified shortcut obtained by replacing the averaging
map in Theorem~\ref{thm:descent} with a norm. For a product
\(\alpha=\prod_i\beta_i\), the norm construction yields
\[
 \alpha^{[M:L]}=\prod_i\Nm_{M/L}(\beta_i).
\]
This is a factorization of a \emph{power} of \(\alpha\), not of
\(\alpha\). Extracting roots can increase degrees, leave \(L\), and
introduce root-of-unity ambiguities. The additive descent proof therefore
does not establish the corresponding global product algorithm. No such
unrestricted product-completeness claim is made for this implementation.

\section{A complete finite-box search for either operation}
\label{sec:bounded}
A complementary approach searches defining polynomials directly. Unlike
inverse coefficient matching over a chosen template, the following
procedure has a simple and explicit bounded completeness guarantee.

\begin{definition}
For an algebraic number \(\beta\), let \(H(\beta)\) be the largest
absolute coefficient of its primitive integer minimal polynomial with
positive leading coefficient. This is a coefficient height, not the
logarithmic Weil height. Fix positive integers \(d,H,k\).
\end{definition}
Enumerate every primitive irreducible integer polynomial
\[
 a_jx^j+\cdots+a_0,\qquad
 1\leq j\leq d,\quad1\leq a_j\leq H,\quad |a_i|\leq H.
\]
Include every exact root of every such polynomial in the catalog.
Nonmonic polynomials must be included: excluding them would discard
nonintegral algebraic numbers and many rational numbers.

For each candidate maximum degree \(e=1,\ldots,d\), use the catalog
entries of degree at most \(e\). For each length \(\ell\leq k\),
enumerate multisets of the first \(\ell-1\) terms or factors. Recover the
last entry exactly as
\[
 \beta_\ell=\alpha-\sum_{i<\ell}\beta_i
 \quad\hbox{or}\quad
 \beta_\ell=\frac{\alpha}{\prod_{i<\ell}\beta_i}.
\]
Zero factors are excluded before division. Test the residual's actual
minimal-polynomial degree and height, then verify the complete identity.
The singleton case uses an empty partial list and simply tests \(\alpha\).

\begin{proposition}[Finite-box completeness]\label{prop:bounded}
If a representation has length at most \(k\), every entry has degree at
most \(d\), and every entry has height at most \(H\), the search above
finds a representation. Searching degree thresholds increasingly finds
the smallest maximum degree among all representations in that box.
\end{proposition}
\begin{proof}
The normalized primitive minimal polynomial of every entry occurs in the
finite enumeration, and the correct conjugate occurs among its roots.
Choose all but one of the entries of a proposed representation; their
multiset is enumerated, and the residual is exactly the remaining entry.
Conversely, every accepted residual is checked against all three bounds,
and the identity is checked exactly. Thus the search is exhaustive in the
specified class and introduces no false positive.
\end{proof}
For the supplied examples one may use
\begin{lstlisting}
BoundedDecompositionSearch[ap, "Product", 3, 1, 2]
BoundedDecompositionSearch[as, "Sum",     3, 1, 2]
\end{lstlisting}
because both cubics have coefficient height one. A degree-three answer then
meets the global lower bound and proves global optimality. This is a second,
independent algorithmic route, although exact residual minimal-polynomial
computations can be substantially more expensive than their modest-looking
catalog sizes suggest.

Increasing the boxes eventually finds every existing finite algebraic
representation. This is a positive semidecision procedure for a fixed
degree threshold: it halts when a representation exists, but an endless
search does not prove nonexistence. Nor does the first result from a growing
box automatically have globally minimum degree. It may precede a
lower-degree representation of larger height or length. A separate lower
bound is needed to turn such a result into a global optimum.

The routine reports \wl{Missing["NoDecompositionWithinBounds", ...]}
only after completing the requested finite box. Resource limits report
\wl{Failure}, not a mathematical nonexistence result. For \(M\) catalog
entries, the number of multisets of \(\ell-1\) entries is
\(\binom{M+\ell-2}{\ell-1}\); even small degree bounds can therefore
produce large searches.

\section{Implementation contracts, usage, and resource limits}
\label{sec:implementation}
\subsection{Public functions}
The package in Appendix~\ref{app:code} is also supplied as the separate
file \wl{RootDecomposition.wl}. Its main entry points are summarized below.

\begin{longtable}{@{}>{\raggedright\arraybackslash}p{0.39\linewidth}>{\raggedright\arraybackslash}p{0.57\linewidth}@{}}
\toprule
\textbf{Function} & \textbf{Contract}\\\midrule
\endhead
\wl{AlgebraicDegree[a]} & Exact absolute degree; rejects inexact or nonalgebraic input.\\[4pt]
\wl{DegreeLowerBound[a]} & Largest prime divisor of the input degree, with rational case one.\\[4pt]
\wl{FieldDecompositionData[a]} & Power-basis coordinates and every subfield of \(\Q(a)\), unless a resource guard stops the calculation.\\[4pt]
\wl{AdditiveDecomposition[a]} & Minimum maximum degree for arbitrary finite sums in \(\Q(a)\).\\[4pt]
\wl{AdditiveDecomposition[a,data]} & Same additive optimization in an explicitly supplied ambient field containing \(a\).\\[4pt]
\wl{GlobalAdditiveDecomposition[a]} & Global additive optimization by normal-closure descent when the smaller field is insufficient.\\[4pt]
\wl{ProductDecomposition[a]} & Complete binary search in \(\Q(a)\), augmented by the independent-family tensor search; global only when separately certified.\\[4pt]
\wl{BoundedDecompositionSearch} & Complete degree--height--length box search for the selected operation.\\[4pt]
\wl{VerifyDecomposition} & Exact identity test for a supplied list and operation.\\[4pt]
\wl{DecompositionExpression} & Displays an inactive sum or product while preserving the separate entries.\\
\bottomrule
\end{longtable}
The resultant helpers \wl{RootSumPolynomial} and
\wl{RootProductPolynomial} implement \eqref{eq:ressum}--\eqref{eq:resprod}.
Their arguments are univariate rational polynomials in a clear symbol.
For interpreting their outputs as monic products over roots, first normalize
the input polynomials to monic form.

A successful result contains
\begin{lstlisting}
<|"Operation" -> "Product", "Terms" -> {...},
  "Degrees" -> {...}, "MaximumDegree" -> ...,
  "LowerBound" -> ..., "IdentityVerified" -> True,
  "GlobalOptimalityCertified" -> ...,
  "OptimalityScope" -> ..., "Method" -> ...|>
\end{lstlisting}
A false \wl{"GlobalOptimalityCertified"} means ``not certified by this
run,'' not ``proved nonoptimal.'' The recorded scope can still establish
an exact internal or bounded minimum. Conversely, when the candidate's
maximum degree equals \wl{"LowerBound"}, the global flag is true even
if the discovery method had a restricted search space. That implication
comes from Theorem~\ref{thm:lpf}, not from search completeness.

\subsection{Reusing an ambient field}
Subfield data can be reused for additive queries in the same field:
\begin{lstlisting}
data = FieldDecompositionData[theta];
r1 = AdditiveDecomposition[alpha1, data];
r2 = AdditiveDecomposition[alpha2, data];
\end{lstlisting}
Each target is checked for membership when coordinates are requested.
A Galois ambient field automatically makes an internally optimal sum
 globally optimal by Theorem~\ref{thm:descent}. Otherwise only an independent
lower bound or a subsequent normal-closure computation supplies that claim.

For the two nonics, inspection is as simple as
\begin{lstlisting}
dp = FieldDecompositionData[ap];
ds = FieldDecompositionData[as];
{Length /@ dp["Subfields"], Length /@ ds["Subfields"]}
(* Mathematical target: {{1,3,3,9}, {1,3,3,9}}. *)
\end{lstlisting}
Data matrices use a fixed power basis. Replacing the primitive element
without converting the matrices changes their meaning; it is not a harmless
presentation change.

\subsection{Resource guards and asymptotic bottlenecks}
Default field limits are \wl{"MaxFieldDegree" -> 64} and
\wl{"MaxSubfields" -> 4096}. The tensor search defaults to
\wl{"MaxTensorFamilies" -> 10000}. The bounded search limits polynomial
enumeration, catalog size, and tuple enumeration separately. Appropriate
limits can be raised or set to \wl{Infinity}, but that does not make a
large problem inexpensive.

The main field bottlenecks are exact factorization over a number field,
rational coefficient growth in the nullspaces, and the number of
intermediate fields. Equation~\eqref{eq:principalmatrix} has \(n\) columns
and at most \(n\deg g\) rows; all entries are rational, but their bit sizes
can grow. Enumerating many canonical intersections can dominate the later
linear membership tests. Partition representations from \cite{szutkoski}
are a natural optimization for large subfield lattices.

The degree guard for a normal closure is checked after a primitive element
has been constructed. Thus it is not a time or memory bound on
\wl{ToNumberField[roots, All]} itself. Users facing large splitting fields
should apply an outer time or memory policy as well. An aborted computation
is an unknown result, never an impossibility certificate.

The reference code favors explicit verification over speed. Potential
improvements include caching coordinate conversions, caching multiplication
matrices, using exact sparse matrices, storing subfields by partitions,
pruning tensor families by field containments, and choosing a smaller
primitive element to reduce coefficient heights. These are engineering
improvements to the algorithms, not missing logical hypotheses in their
stated mathematical guarantees.

\subsection{Output quality versus degree optimality}
Even a globally degree-optimal decomposition need not have small
coefficients or resemble the particular cubics written in the question.
The additive trace normalization and checked rational product normalization
make the examples more readable, but do not solve a secondary height
minimization problem. Additional criteria could be optimized only after
fixing the optimal maximum degree. Examples include the number of terms,
coefficient height, number of real factors, or expression size.

We allow complex algebraic entries even when the input is real. Restricting
all entries to be real defines a different problem and is not imposed by
the global algorithms. The two main examples and the displayed auxiliary
examples already have real optimal entries, so this distinction does not
change their conclusions.

\section{Executed verification and reproducibility}\label{sec:verification}
\subsection{What was actually run}
The accompanying \wl{verify\_exact.py} was executed using SymPy 1.14.0.
It implements the principal-field nullspaces and intersections independently,
using exact algebraic-field arithmetic and rational matrices. It is not a
wrapper around the Wolfram package. The run completed with all assertions
passing; its recorded elapsed time was 7.82 seconds in the preparation
environment. This is a SymPy verification time, not a Mathematica benchmark.

The machine-readable report \wl{verification\_results.json} records:
\begin{center}
\begin{tabular}{@{}p{0.51\linewidth}p{0.43\linewidth}@{}}
\toprule
\textbf{Check} & \textbf{Executed result}\\\midrule
Both nonic resultant identities & Exact polynomial equality\\
Nonic irreducibility and real branches & Degree nine; one real root each\\
Recovery certificates & All six polynomial remainders zero\\
Two nonic subfield lattices & Degrees \(1,3,3,9\) for each\\
Binary product and additive searches & Maximum degree three\\
Three-term quadratic sum & Degree bound two; three terms\\
Three-factor quadratic product & Degree bound two; three factors\\
Degree-eight field lattice & Sixteen subfields\\
Both external-factor sextic lattices & Degrees \(1,3,6\)\\
\bottomrule
\end{tabular}
\end{center}
Every computed subfield basis is additionally checked for containing one
and for closure under multiplication of its basis elements. This is a useful
independent check that mere linear subspaces have not been confused with
subfields.

\subsection{What was not run}
The available Wolfram connector was actually invoked, but returned HTTP 404
before any kernel evaluation. Therefore the package, its local regression
suite, and the optional degree-24 normal-closure example were not executed
in Mathematica. The source was statically reviewed, including balanced
brackets, nested comments, and strings. That check does not validate all
Wolfram evaluation semantics.

The Python verification does not test the Wolfram normal-closure
construction or reproduce a Mathematica transcript. The proof of
Theorem~\ref{thm:descent} establishes the mathematical global additive
algorithm; runtime behavior of its Wolfram realization remains for the
supplied local suite and additional tests to assess.

\subsection{Reproduction commands}
With SymPy installed, run
\begin{lstlisting}[language={}]
python verify_exact.py
\end{lstlisting}
This regenerates the JSON report, including a new measured elapsed time.
In Mathematica, execute
\begin{lstlisting}
TestReport["/your/path/tests.wlt"]
\end{lstlisting}
The test file loads the package relative to itself. It includes exact
identities, resultants, degree and branch tests, all recovery congruences,
subfield recovery, rational and zero cases, rejection of inexact input, the
three-term and three-factor examples, and the internal/global distinction.
The more expensive full normal-closure example is available separately in
\wl{examples.wl} as an optional commented section.

The PDF is reproducible with two \wl{pdflatex} passes; the script
\wl{build.sh} supplies the commands.

Keep the source file \wl{RootDecomposition.wl} beside the \LaTeX{}
source, because the full implementation below is included directly from
that file. This prevents the printed source and executable source from
drifting apart.

\section{Conclusions}
Both supplied degree-nine numbers have globally optimal decompositions
into the requested cubic roots. Exact polynomial congruences certify their
identities, and the largest-prime-factor obstruction forces the optimum
three even for arbitrarily long sums or products outside the input field.

Subfield recovery provides a template-free route to those answers. Rational
linear algebra handles all finite sums in a fixed field; normal-closure
descent makes additive minimization globally complete. Binary products also
have an exact linear test, while independent multi-factor families admit a
rank-one tensor certificate. Finite-box enumeration gives another exact,
but bounded, discovery procedure.

The product guarantees must remain qualified. Three quadratic factors can
outperform every quadratic pair, and the sextic counterexamples require
leaving the input field to attain their optimum. The returned certificates
therefore state the scope of each optimality claim.

\clearpage
\appendix
\section{Full Wolfram Language reference implementation}\label{app:code}
This is the exact companion source file. Execution status and mathematical
scope are documented in Section~\ref{sec:verification}; the listing is not
presented as runtime-tested Wolfram output. The independent verification and
the local \wl{tests.wlt} suite are separate files in the archive.
\lstinputlisting[numbers=left,numberstyle=\tiny\color{commentgray},numbersep=7pt,
 xleftmargin=16pt,basicstyle=\ttfamily\scriptsize]{RootDecomposition.wl}

\vspace{1em}
\begin{thebibliography}{99}
\small\setlength{\itemsep}{3pt}\setlength{\parskip}{0pt}
\addcontentsline{toc}{section}{References}
\bibitem{sequestion}
Vladimir Reshetnikov, Mathematica StackExchange question 105933, concerning
products and sums of lower-degree algebraic roots; answers by yarchik and
Daniel Lichtblau (2019).
\url{https://mathematica.stackexchange.com/q/105933/7288}.
The question's two examples are reproduced from the user's supplied text.

\bibitem{szutkoski}
J. Szutkoski and M. van Hoeij,
\emph{The Complexity of Computing all Subfields of an Algebraic Number
Field}, arXiv:1606.01140v3, 20 November 2017.
\url{https://arxiv.org/html/1606.01140v3}.

\bibitem{milne}
J. S. Milne, \emph{Fields and Galois Theory}, version 5.10, September 2022.
In particular Chapters 3--5 on the Galois correspondence, Galois groups,
primitive elements, norms, and traces.
\url{https://www.jmilne.org/math/CourseNotes/FT.pdf}.

\bibitem{root}
Wolfram Research, \emph{Root}.
\url{https://reference.wolfram.com/language/ref/Root.html}.

\bibitem{rootreduce}
Wolfram Research, \emph{RootReduce}.
\url{https://reference.wolfram.com/language/ref/RootReduce.html}.

\bibitem{minpoly}
Wolfram Research, \emph{MinimalPolynomial}.
\url{https://reference.wolfram.com/language/ref/MinimalPolynomial.html}.

\bibitem{tonumberfield}
Wolfram Research, \emph{ToNumberField}.
\url{https://reference.wolfram.com/language/ref/ToNumberField.html}.

\bibitem{anpoly}
Wolfram Research, \emph{AlgebraicNumberPolynomial}.
\url{https://reference.wolfram.com/language/ref/AlgebraicNumberPolynomial.html}.

\bibitem{factorlist}
Wolfram Research, \emph{FactorList}.
\url{https://reference.wolfram.com/language/ref/FactorList.html}.

\bibitem{polyrem}
Wolfram Research, \emph{PolynomialRemainder}.
\url{https://reference.wolfram.com/language/ref/PolynomialRemainder.html}.

\bibitem{resultant}
Wolfram Research, \emph{Resultant}.
\url{https://reference.wolfram.com/language/ref/Resultant.html}.

\bibitem{linearsolve}
Wolfram Research, \emph{LinearSolve}.
\url{https://reference.wolfram.com/language/ref/LinearSolve.html}.

\bibitem{countroots}
Wolfram Research, \emph{CountRoots}.
\url{https://reference.wolfram.com/language/ref/CountRoots.html}.

\bibitem{irreducible}
Wolfram Research, \emph{IrreduciblePolynomialQ}.
\url{https://reference.wolfram.com/language/ref/IrreduciblePolynomialQ.html}.
\end{thebibliography}
{\small Official Wolfram documentation consulted 6 September 2026.}
\end{document}
